
Consider a system with dynamics described by

\begin{equation}
    \mathbf{x}(k) = \mathbf{A} x(k) + \mathbf{B} \mathbf{u}(k)
\end{equation}

The objective is to drive the state $x(K)$ to the origin from a given initial condition while respecting
state and control constraints in the form

\begin{equation}
	\mathbf{x}(k) \in \mathcal{X}
\end{equation}
	
\begin{equation}
	\mathbf{u}(k) \in \mathcal{U}
\end{equation}


\begin{equation}
	\mathcal{U} = \left\{
		\mathbf{u} \in \mathbb{R}^p : S_u u \leq b_u	
	\right\}
\end{equation}

\noindent
and

\begin{equation}
	\mathcal{X} = \left\{
		x \in \mathbb{R}^n : S_x x \leq b_x
	\right\}
\end{equation}


\section{Cost Function}

The solution, if it exists, will not be unique. Therefore, it is necessary to adopt a criterion that allows for choosing the best sequence of control 
actions to perform the task.


The control problem consists of bringing the states to the origin, so we will employ a cost in the form

\begin{equation}
	J = \blue{\Vert x(N) \Vert^2_{P_f}} + \orange{\sum^{N-1}_{k=0} \Big[ \Vert x(k) \Vert^2_Q +  \Vert u(k) \Vert^2_R \Big]}
\end{equation}

The expression for $J$ consists of a \blue{terminal cost} term and a trajectory \orange{cost term}. 

With $P_f$, $Q$, $R$ being symmetric and positive-definite, with $P_f$ chosen so that the terminal cost term $\norma{\hat{x}(k + N|k) - \bar{x}}{P_f}$ corresponds 
to the value function of the infinite horizon Discrete-Time Linear Quadratic Regulator (DLQR).

The cost function can be rewritten using the notation:

\begin{equation}
	\mathbf{u} = 
	\left[
		\begin{matrix}
			u(0) \\		
			u(1) \\		
			\vdots \\
			u(N-1)
		\end{matrix}
	\right]
	\in \mathbb{R}^{pN},
	\mathbf{x} = 
	\left[
		\begin{matrix}
			x(1) \\ 
			x(2) \\ 
			\vdots \\
			x(N)
		\end{matrix}
	\right]
	\in \mathbb{R}^{nN}
\end{equation}

\noindent
and then the sum equation

\begin{equation}
	\sum^{N-1}_{k=0} \Vert u(k) \Vert^2_R = u^T(0) R u(0) + u^t(1)Ru(1) + \dotsm + u^T (N-1) R u(N-1)
\end{equation}

\noindent
becomes

\begin{equation}
	\underbrace{
		\left[
			\begin{matrix}
				u^T(0) & u^T(1) & \dotsi & u^T(N - 1) 
			\end{matrix}
		\right]
	}_{\mathbf{u}^T}
	\left[
		\begin{matrix}
			R & 0 & \dotsi & 0 \\
			0 & R & \dotsi & 0 \\
			\vdots & \vdots & \ddots & \vdots \\
			0 & 0 & \dotsi & R \\
		\end{matrix}
	\right]
	\underbrace{
		\left[
			\begin{matrix}
				u(0) \\
				u(1) \\
				\vdots \\
				u(N - 1)
			\end{matrix}
		\right]
	}_{\mathbf{u}}
\end{equation}

Using the notation

\begin{equation}
	\mathbf{R} = \left[
		\begin{matrix}
			R & 0 & \dotsi & 0 \\
			0 & R & \dotsi & 0 \\
			\vdots & \vdots & \ddots & \vdots \\
			0 & 0 & \dotsi & R 
		\end{matrix}
	\right] = diag_N(R)
\end{equation}

\noindent
we get to 

\begin{equation}
	\sum_{k=0}^{N-1} \norma{u(k)}{R} = \norma{\mathbf{u}}{\mathbf{R}}
\end{equation}

In a similar way, we can write

\begin{equation}
	\begin{gathered}
		\norma{x(N)}{P_f} + \sum_{k=0}^{N-1} \norma{x(k)}{Q} = x^T(0) Q x(0) \\
		+ x^T(1)Qx(1) + \dotsi + x^T(N-1)Qx(N-1) + x^T(N)P_fx(N) \\
		= x^T(0)Qx(0) +  \\
		\underbrace{
			\left[
				\begin{matrix}
					x^T(1) & x^T(2) & \dotsi & x^T(N)		
				\end{matrix}
			\right]
		}_{\mathbf{x}^T}
		\left[
			\begin{matrix}
				Q & 0 & \dotsi & 0 \\
				0 & Q & \dotsi & 0 \\
				\vdots & \vdots & \ddots & \vdots \\
				0 & 0 & \dotsi & P_f
			\end{matrix}
		\right]
		\underbrace{
			\left[
				\begin{matrix}
					x(1) \\
					x(2) \\
					\vdots \\
					x(N)
				\end{matrix}
			\right]
		}_{\mathbf{x}}
	\end{gathered}
\end{equation}

\noindent
with the notation

\begin{equation}
	Q = 
	\left[
		\begin{matrix}
			Q & 0 & \dotsi & 0 \\
			0 & Q & \dotsi & 0 \\
			\vdots & \vdots & \ddots & \vdots \\
			0 & 0 & \dotsi & P_f
		\end{matrix}
	\right]
	= 
	\left[
		\begin{matrix}
			diag_{N-1}(Q) & 0 \\
			0 & P_f
		\end{matrix}
	\right]
	\in \mathbb{R}^{nN \times nN}
\end{equation}

\noindent
as

\noindent
\begin{equation}
	\norma{x(N)}{{_f}} + \sum^{N - 1}_{k=0} \norma{x(k)}{Q} = \norma{x(0)}{Q} + \norma{\mathbf{x}}{\mathbf{Q}} 
\end{equation}

So, the cost function can be written in the form

\noindent
\begin{equation}
	J(\mathbf{x}, \mathbf{u}) = \norma{x(0)}{Q} + \norma{\mathbf{x}}{\mathbf{Q}} + \norma{\mathbf{u}}{\mathbf{R}}
\end{equation}